\documentclass[12pt]{article}
\usepackage{amsmath, amsthm, amsfonts, mathtools, xcolor, caption}

\numberwithin{equation}{section}

\usepackage[style=numeric,sorting=none,giveninits=true]{biblatex}
\addbibresource{bibliography.bib}

\usepackage{etoolbox}
\usepackage[margin=1in]{geometry}

\DeclareMathAlphabet{\mymathbb}{U}{BOONDOX-ds}{m}{n}

\usepackage{comment}

%%%%%%%%%%%%%%%%%%%%%%%%%%%%%%%%%%%%%%%%%
%       Theorem Environments            %
%%%%%%%%%%%%%%%%%%%%%%%%%%%%%%%%%%%%%%%%%

\newtheorem{theorem}{Theorem}[section]
\newtheorem{lemma}[theorem]{Lemma}
\newtheorem{corollary}[theorem]{Corollary}
\newtheorem{proposition}[theorem]{Proposition}
\theoremstyle{definition}
\newtheorem{remark}{Remark}
\newtheorem{definition}{Definition}
\newtheorem{notation}{Notation}
\newtheorem{example}{Example}
\newtheorem{conjecture}[theorem]{Conjecture}
\newtheorem{openproblem}[theorem]{Open Problem}

%%%%%%%%%%%%%%%%%%%%%%%%%%%%%%%%%%%%%%%%%
%               New Commands            %
%%%%%%%%%%%%%%%%%%%%%%%%%%%%%%%%%%%%%%%%%

\newcommand{\ev}[1]{( #1 )} 
\newcommand{\ew}[1]{\left( #1 \right)} 

\newcommand{\N}{\mathbb{N}}
\newcommand{\Z}{\mathbb{Z}}
\newcommand{\C}{\mathbb{C}}
\newcommand{\Q}{\mathbb{Q}}
\newcommand{\R}{\mathbb{R}}

%%%%%%%%%%%%%%%%%%%%%%%%%%%%%%%%%%%%%%%%%
%               Document                %
%%%%%%%%%%%%%%%%%%%%%%%%%%%%%%%%%%%%%%%%%

\let\underbrace\LaTeXunderbrace
\let\overbrace\LaTeXoverbrace

\usepackage{parskip}

\begin{document}

\section*{Challenge 9: Union-Closed Sets Conjecture}
\refstepcounter{section}

\begin{definition}[Union-closed set]
A set $\mathcal{F}$ of sets is \textit{union-closed} if for any two sets in $\mathcal{F}$, their union is also in $\mathcal{F}$. 
For a nonempty finite set $\mathcal{F}$, the \emph{density} of an element $x$ with respect to $\mathcal{F}$ is defined
to be 
\[
\frac{|\{ F \in \mathcal{F} : x \in F \}|}{|\mathcal{F}|}
\]
\end{definition}

\begin{conjecture}[Union-closed sets conjecture \mbox{\cite{nagel2023}}]
For every finite union-closed set $\mathcal{F}\notin\{\emptyset, \{\emptyset\}\}$, there exists an element $x$ of density at least one half.
\end{conjecture}

We are interested in the following parametrized equialent version of the conjecture.
\begin{conjecture}[Scaled union-closed sets conjecture]
Let $r\in \mathbb{N}$.
For every finite union-closed set $\mathcal{F}\notin\{\emptyset, \{\emptyset\}\}$, there exists an element $x$ of density at least $\frac{1}{2}-\frac{1}{r+2}$.
\end{conjecture}

Here we are interested in proving this conjecture when specialised to particular values of $r$. 
For $r=0$ the bound equals $0$ and the statement is trivial.
The best published lower bound on the density is $\approx 0.382$. It was obtained via the entropy method by Chase and Lovett \cite{chase2022approximate} and independently by Sawin~\cite{sawin2022unionclosed}. This settles the cases $r\leq 6$; all values $r\geq 7$ remain open to date.

%%%%%%%%%%%%%%%%%%%%%%%%%%%%%%%%%%%%%%%%%%%%%%%%%%
\printbibliography
%%%%%%%%%%%%%%%%%%%%%%%%%%%%%%%%%%%%%%%%%%%%%%%%%%
\end{document}
